% Part: first-order-logic
% Chapter: completeness
% Section: completeness-theorem

\documentclass[../../../include/open-logic-section]{subfiles}

\begin{document}

\iftag{FOL}
      {\olfileid{fol}{com}{cth}}
      {\olfileid{pl}{com}{cth}}

\olsection{Teorema Kelengkapan}

\begin{explain}
Kanthi nggabungake asil-asil iki, kita oleh teorema kelengkapan.
\end{explain}

\begin{thm}[Teorema Kelengkapan]
\ollabel{thm:completeness}
Upamana $\Gamma$ minangka himpunan !!{sentence}s. Yen $\Gamma$
konsisten, himpunan kasebut bisa dicukupi.
\end{thm}

\begin{proof}
Upamana $\Gamma$ konsisten.\iftag{FOL}{ Miturut
  \olref[hen]{lem:henkin}, ana himpunan konsisten kang jenuh
  $\Gamma' \supseteq \Gamma$.}{} Miturut \olref[lin]{lem:lindenbaum},
ana $\Gamma^* \supseteq \iftag{FOL}{\Gamma'}{\Gamma}$ kang konsisten
lan !!{complete}.
\iftag{FOL}{
  Amarga $\Gamma' \subseteq \Gamma^*$, kanggo saben
  !!{formula}~$!A(x)$, $\Gamma^*$ ngemot !!a{sentence} kanthi wujud
  \iftag{prvEx}
        {$\lexists[x][!A(x)] \lif !A(c)$}
        {$\lnot\lforall[x][!A(x)] \lif \lnot !A(c)$}
  lan mula $\Gamma^*$ jenuh. Yen $\Gamma$ ora ngemot~$\eq$, miturut}{Miturut}
\olref[mod]{lem:truth},
$\iftag{FOL}{\Sat{M(\Gamma^*)}}{\pSat{v(\Gamma^*)}}{!A}$ yen lan mung
yen $!A \in \Gamma^*$. Saka kene mligine nderek yen kanggo kabeh
$!A \in \Gamma$, $\iftag{FOL}{\Sat{M(\Gamma^*)}}{\pSat{v(\Gamma^*)}}{!A}$,
mula $\Gamma$ bisa dicukupi.\iftag{FOL}{
  Yen $\Gamma$ ngemot~$\eq$, miturut \olref[ide]{lem:truth}, kanggo
  kabeh !!{sentence}s~$!A$, $\Sat{\equivclass{M}{\approx}}{!A}$ yen lan
  mung yen $!A \in \Gamma^*$. Mligine,
  $\Sat{\equivclass{M}{\approx}}{!A}$ kanggo kabeh $!A \in \Gamma$,
  mula $\Gamma$ bisa dicukupi.}{}
\end{proof}

\begin{cor}[Teorema Kelengkapan, Versi Kapindho]
\ollabel{cor:completeness}
Kanggo kabeh $\Gamma$ lan !!{sentence}s~$!A$: yen $\Gamma \Entails !A$
mula $\Gamma \Proves !A$.
\end{cor}

\begin{proof}
Elinga yen $\Gamma$ ing \olref{cor:completeness} lan
\olref{thm:completeness} dikuwantifikasi universal. Supaya ora kliru,
pratelan \olref{thm:completeness} kita tulis maneh nganggo variabel
liyane: kanggo saben himpunan !!{sentence}s~$\Delta$, yen $\Delta$
konsisten, himpunan kasebut bisa dicukupi. Miturut kontraposisi, yen
$\Delta$ ora bisa dicukupi, mula $\Delta$ ora konsisten. Iki bakal
digunakake kanggo mbuktekake korolari kasebut.

Upamana $\Gamma \Entails !A$. Banjur $\Gamma \cup \{\lnot !A\}$ ora
bisa dicukupi miturut \olref[syn][sem]{prop:entails-unsat}. Kanthi
njupuk $\Gamma \cup \{\lnot !A\}$ minangka $\Delta$, versi sadurunge
saka \olref{thm:completeness} menehi yen $\Gamma \cup \{\lnot !A\}$
ora konsisten. Miturut
\iftag{FOL}{%
  \tagrefs{prfAX/{fol:axd:prv:prop:prov-incons},
    prfSC/{fol:seq:prv:prop:prov-incons},
    prfND/{fol:ntd:prv:prop:prov-incons},
    prfTab/{fol:tab:prv:prop:prov-incons}}}{%
  \tagrefs{prfAX/{pl:axd:prv:prop:prov-incons},
    prfSC/{pl:seq:prv:prop:prov-incons},
    prfND/{pl:ntd:prv:prop:prov-incons},
    prfTab/{pl:tab:prv:prop:prov-incons}}},
$\Gamma \Proves !A$.
\end{proof}

\tagprob{FOL}
\begin{prob}
Gunakna \olref[fol][com][cth]{cor:completeness} kanggo mbuktekake
\olref[fol][com][cth]{thm:completeness}, lan kanthi mangkono tuduhna
yen rong formulasi teorema kelengkapan kasebut ekuivalen.
\end{prob}
\tagendprob

\tagprob{notFOL}
\begin{prob}
Gunakna \olref[pl][com][cth]{cor:completeness} kanggo mbuktekake
\olref[pl][com][cth]{thm:completeness}, lan kanthi mangkono tuduhna
yen rong formulasi teorema kelengkapan kasebut ekuivalen.
\end{prob}
\tagendprob

\tagprob{FOL}
\begin{prob}
Supaya sawijining sistem !!a{derivation} jangkep, aturan-aturane kudu cukup kuwat
kanggo mbuktekake yen saben himpunan kang ora bisa dicukupi iku ora
konsisten. Aturan !!{derivation} endi wae kang dibutuhake kanggo
mbuktekake kelengkapan? Apa ana aturan kasebut kang ora digunakake ing
bukti iki? Kanggo mangsuli pitakon-pitakon iki, gawenen dhaptar utawa
diagram kang nuduhake aturan !!{derivation} endi wae kang digunakake ing
asil-asil kang tumuju marang bukti
\olref[fol][com][cth]{thm:completeness}. Catheten uga saben panganggone
aturan kang ora kasebut kanthi cetha ing bukti-bukti kasebut.
\end{prob}
\tagendprob

\tagprob{notFOL}
\begin{prob}
Supaya sawijining sistem !!a{derivation} jangkep, aturan-aturane kudu cukup kuwat
kanggo mbuktekake yen saben himpunan kang ora bisa dicukupi iku ora
konsisten. Aturan !!{derivation} endi wae kang dibutuhake kanggo
mbuktekake kelengkapan? Apa ana aturan kasebut kang ora digunakake ing
bukti iki? Kanggo mangsuli pitakon-pitakon iki, gawenen dhaptar utawa
diagram kang nuduhake aturan !!{derivation} endi wae kang digunakake ing
asil-asil kang tumuju marang bukti
\olref[pl][com][cth]{thm:completeness}. Catheten uga saben panganggone
aturan kang ora kasebut kanthi cetha ing bukti-bukti kasebut.
\end{prob}
\tagendprob

\end{document}
